\documentclass[Riemann_Hypothesis.tex]{subfiles}
\begin{document}


\pagebreak
\color{red}
\part{Resolution of the conjecture}
\color{black}

\color{red}
\section{First definitions and Weil space}
\color{black}

\begin{DefNefEn}[label=df:EspaceSchwartz]{Schwartz space}
    Let $\mathbb{R}$ be the field of real numbers. \par
    Following the work in functional analysis introduced by Laurent Schwartz 
    (see~\cite{schwartz1966_distributions_en}), we call Schwartz space, denoted by $\mathcal{S}(\mathbb{R})$, the vector space of functions $f \in C^\infty(\mathbb{R}, \mathbb{C})$ that are rapidly decreasing. \par
    A function $f$ belongs to the space $\mathcal{S}(\mathbb{R})$ if and only if, for every pair of integers $(\alpha, \beta) \in \mathbb{N}^2$, we have:
    \begin{equation*}
        \sup_{x \in \mathbb{R}} \left| x^\alpha f^{(\beta)}(x) \right| \;<\; +\infty
    \end{equation*}
    where $f^{(\beta)}$ denotes the derivative of order $\beta$ of the function $f$.
\end{DefNefEn}


\begin{DefNefEn}[label=df:EspaceWeil_fr]{Weil test function space}
    Let $\mathbb{R}_+^*$ be the multiplicative group of strictly positive real numbers. \par
    According to the work of André Weil~\cite{weil1952_formules_en}, 
    we call the Weil test function space, denoted by $\mathcal{W}$, the set of functions $f \in C^\infty(\mathbb{R}_+^*, \mathbb{C})$ such that the function $g : t \mapsto f(e^t)$ belongs to the Schwartz space $\mathcal{S}(\mathbb{R})$. \par
    In order to formulate an equivalent condition, we introduce the differential operator $D$ defined on $C^\infty(\mathbb{R}_+^*, \mathbb{C})$ by:
    \begin{equation*}
        D(f)(x) = x f'(x)
    \end{equation*}
    We will equivalently say that $f \in \mathcal{W}$ if and only if, for any pair of integers $(k, m) \in \mathbb{N}^2$, we have:
    \begin{equation*}
        \sup_{x \in \mathbb{R}_+^*} \left| (\ln x)^k D^m(f)(x) \right| \;<\; +\infty
    \end{equation*}
    where $D^m$ denotes the $m$-th iteration of the operator $D$.
\end{DefNefEn}



\begin{DefNefEn}[label=df:MesureHaar]{Multiplicative Haar measure}
    Following the work of Alfred Haar (see Bourbaki~\cite{Haar_bourbaki_integration_chap7_en}), on any locally compact topological group, there exists a non-trivial regular Radon measure, unique up to a multiplicative constant, which is translation-invariant. \par
    Let $G = \left(\mathbb{R}_+^*, \times\right)$ be the multiplicative topological group of strictly positive real numbers. \par
    We call the \textbf{Haar measure} of the group $G$, which we will denote by $d^\times x$, the unique measure (up to a scaling factor) invariant under the action of multiplication. It is defined by:
    \begin{equation*}
        d^\times x \;=\; \frac{\mathrm{d}x}{x}
    \end{equation*}
    where $\mathrm{d}x$ denotes the usual Lebesgue measure on $\mathbb{R}$. \par
    Indeed, let $a$ be a strictly positive real number. \par
    Consider the left translation map on the group, defined by the multiplication $x \mapsto ax$. \par
    By differentiating this map, we have $\mathrm{d}(ax) = a \, \mathrm{d}x$. \par
    Evaluating the measure on the translated element gives:
    \begin{equation*}
        d^\times (ax) \;=\; \frac{\mathrm{d}(ax)}{ax} \;=\; \frac{a \, \mathrm{d}x}{ax} \;=\; \frac{\mathrm{d}x}{x}
    \end{equation*}
    Since $d^\times (ax) = d^\times x$, we deduce that the measure thus constructed is indeed invariant under dilation. 

    It is this invariance that ensures the structural compatibility with the Mellin transform and the Weil Hermitian form.
\end{DefNefEn}


\begin{lemmaEn}[label=lm:DensiteWeil]{Density of the Weil space}
    Let $L^2(\mathbb{R}_+^*, d^\times x)$ be the Hilbert space of complex-valued square-integrable functions with respect to the Haar measure. \par
    The Weil test function space $\mathcal{W}$ is dense in $L^2(\mathbb{R}_+^*, d^\times x)$.
\end{lemmaEn}

\begin{proof}
    \begin{part_proof}
        Let $\Phi : \mathbb{R} \to \mathbb{R}_+^*$ be the map defined by $\Phi(t) = e^t$. \par
        Since the exponential function is a diffeomorphism from $\mathbb{R}$ onto $\mathbb{R}_+^*$, let $J$ be the composition operator which to any function $f \in L^2(\mathbb{R}_+^*, d^\times x)$ associates the function $J(f) = f \circ \Phi$. \par
        We evaluate the squared norm of $J(f)$ in the usual space $L^2(\mathbb{R}, dt)$:
        \begin{equation*}
            \int_{\mathbb{R}} \left| J(f)(t) \right|^2 \, dt \;=\; \int_{\mathbb{R}} \left| f(e^t) \right|^2 \, dt
        \end{equation*}
        We perform the change of variable $x = e^t$. \par
        Since $dx = e^t \, dt = x \, dt$, we deduce that $dt = \frac{dx}{x} = d^\times x$. \par
        Applying this change of variable to the bounds and the integration measure, we obtain:
        \begin{equation*}
            \int_{\mathbb{R}} \left| f(e^t) \right|^2 \, dt \;=\; \int_{0}^{+\infty} \left| f(x) \right|^2 \, \frac{dx}{x} \;=\; \int_{\mathbb{R}_+^*} \left| f(x) \right|^2 \, d^\times x
        \end{equation*}
        Since the norms are preserved, we deduce that the operator $J$ is a surjective isometry from the Hilbert space $L^2(\mathbb{R}_+^*, d^\times x)$ to the Hilbert space $L^2(\mathbb{R}, dt)$. \medskip\par
        By definition of the space $\mathcal{W}$ (see {Definition}~\ref{df:EspaceWeil}), the image of $\mathcal{W}$ under the isomorphism $J$ is exactly the Schwartz space $\mathcal{S}(\mathbb{R})$. \par
        According to the fundamental theorems of functional analysis 
        (see the work of Laurent Schwartz~\cite{schwartz1966_distributions_en}), the space $\mathcal{S}(\mathbb{R})$ is dense in $L^2(\mathbb{R}, dt)$. \par
        Since $J(\mathcal{W})$ is dense in $L^2(\mathbb{R}, dt)$ and $J^{-1}$ is a continuous map (since $J$ is an isometry), we deduce that $\mathcal{W} = J^{-1}(\mathcal{S}(\mathbb{R}))$ is dense in $L^2(\mathbb{R}_+^*, d^\times x)$.
    \end{part_proof}
\end{proof}

\begin{DefNefEn}[label=df:TransformeeMellin]{Mellin transform}
    Let $f$ be a complex-valued function, defined and locally integrable on the multiplicative group $\mathbb{R}_+^*$. \par
    Following the work of Robert Hjalmar Mellin 
    (see the work of Jean Dieudonné~\cite{Mellin_dieudonne_calcul_infinitesimal_en}), 
    we call the Mellin transform of $f$, denoted by $\mathcal{M}[f]$, the function of 
    the complex variable $s$ defined by:
    \begin{equation*}
        \mathcal{M}[f](s) \;=\; \int_{0}^{+\infty} f(x) x^{s-1} \, dx \;=\; \int_{0}^{+\infty} f(x) x^s \, \frac{dx}{x}
    \end{equation*}
    for any $s \in \mathbb{C}$ such that this integral is absolutely convergent. \par
    Let $\mathcal{D}_f \subset \mathbb{C}$ be the domain of absolute convergence of this integral, which usually constitutes a vertical strip in the complex plane.
\end{DefNefEn}


\begin{propositionEn}[label=pr:MellinHolomorpheWeil]{Holomorphy of the Mellin transform on the Weil space}
    For any function $f \in \mathcal{W}$, its Mellin transform 
    $\mathcal{M}[f]$ is an entire function, and therefore 
    holomorphic on the whole complex plane $\mathbb{C}$.
\end{propositionEn}

\begin{proof}
    \begin{part_proof}
        Let $f \in \mathcal{W}$. \par
        Let $s \in \mathbb{C}$. \par
        Let $\sigma = \Re(s)$ be the real part of the complex number $s$. \par
        According to the definition of the Weil space (see Definition \ref{df:EspaceWeil}), we know that the function $g : t \mapsto f(e^t)$ belongs to the Schwartz space $\mathcal{S}(\mathbb{R})$. \par
        We deduce that for any integer $N \in \mathbb{N}$, there exists a constant $C_N > 0$ such that for all $t \in \mathbb{R}$:
        \begin{equation*}
            |g(t)| \;=\; |f(e^t)| \;\leq\; \frac{C_N}{1 + e^{N|t|}}
        \end{equation*}
        By making the change of variable $x = e^t$, since $x \in \mathbb{R}_+^*$, we deduce that for any integer $N \in \mathbb{N}$, there exists a real number $M_N > 0$ such that:
        \begin{equation*}
            |f(x)| \;\leq\; M_N \min\left(x^N, x^{-N}\right)
        \end{equation*}
        This upper bound proves that the function $f$ decays faster than any power of $x$ as $x \to 0$, and faster than any power of the inverse of $x$ as $x \to +\infty$. \par
        We evaluate the modulus of the integrand defining the Mellin transform:
        \begin{equation*}
            \left| f(x) x^{s-1} \right| \;=\; |f(x)| x^{\sigma-1}
        \end{equation*}
        Since we have the previous upper bound for any integer $N$, we can choose 
        $N > |\sigma - 1| + 1$, and we deduce that the function 
        $x \mapsto |f(x)| x^{\sigma-1}$ is integrable on $\mathbb{R}_+^*$. \medskip\par

        Thus, for any $s \in \mathbb{C}$, the integral defining 
        $\mathcal{M}[f](s)$ is absolutely convergent. \par
        Moreover, the function $s \mapsto f(x) x^{s-1}$ is entire 
        for all $x \in \mathbb{R}_+^*$. \par
        By the holomorphy theorem for integrals with parameters, 
        based on Henri Lebesgue's dominated convergence theorem 
        (see~\cite{lebesgue1902_these_en}), we can locally dominate the integrand 
        independently of $s$ on any compact set of $\mathbb{C}$.
    \end{part_proof}


    \begin{part_proof}
        According to the holomorphy theorem for integrals depending on a parameter, 
        based on Henri Lebesgue's measure theory (see\cite{lebesgue1902_these_en}), for a function defined by an integral to be holomorphic on an open set $U$ of $\mathbb{C}$, it is sufficient that the integrand be measurable with respect to the integration variable, holomorphic with respect to the parameter on $U$, and that it admits a domination by an integrable function independent of the parameter on any compact set included in $U$. \par
        Let $s_0 \in \mathbb{C}$. \par
        Let $K$ be a compact neighborhood of $s_0$ in the complex plane $\mathbb{C}$, such as a closed disk centered at $s_0$ of strictly positive radius. \par
        Since we have previously established the existence of an integrable function $\varphi$ on $\mathbb{R}_+^*$ such that for all $s \in K$ and for all $x \in \mathbb{R}_+^*$, we have $\left| f(x) x^{s-1} \right| \leq \varphi(x)$, we deduce that the hypotheses of the holomorphy theorem are satisfied on the interior of this compact set $K$. \par
        We deduce that the function $\mathcal{M}[f]$ is holomorphic on the interior of $K$. \par
        Since the point $s_0$ belongs to the interior of $K$, we deduce that the function $\mathcal{M}[f]$ is complex-differentiable at the point $s_0$. \par
        Since we have this result for any point $s_0 \in \mathbb{C}$, we deduce that the function $\mathcal{M}[f]$ is complex-differentiable at every point in the complex plane. \par
        By definition, we deduce that the function $\mathcal{M}[f]$ is holomorphic on the whole complex plane $\mathbb{C}$, that is to say, it is an entire function.
    \end{part_proof}

\end{proof}

\color{red}
\section{Weil Hermitian form and Weil positivity criterion}
\color{black}


\begin{DefNefEn}[label=df:FormeHermitienneWeil]{Weil Hermitian form}
    Let $\mathcal{W}$ be the space of Weil test functions on the multiplicative group $\left(\mathbb{R}_+^*, \times\right)$. \par
    Let $f$ and $g$ be two functions in $\mathcal{W}$. \par
    Let $g^*$ be the involutive function defined for all $x \in \mathbb{R}_+^*$ by:
    \begin{equation*}
        g^*(x) \;=\; \frac{1}{x} \, \overline{g\left(\frac{1}{x}\right)}
    \end{equation*}
    Let $\Phi = f * g^*$ denote the multiplicative convolution product of $f$ and $g^*$. \par
    According to André Weil's explicit formulas, we call the \textbf{Weil Hermitian form}, which we will denote by $\langle f, g \rangle_{\mathcal{HP}}$, the sesquilinear form defined by the action of the geometric distribution on the test function $\Phi$:
    \begin{equation*}
        \langle f, g \rangle_{\mathcal{HP}} \;:=\; W_{\text{pol}}(\Phi) - W_{\text{arith}}(\Phi) - W_{\text{arch}}(\Phi)
    \end{equation*}
    where the different functionals acting on $\Phi$ are explicitly given analytically by:
    \begin{align*}
        W_{\text{pol}}(\Phi) \;&=\; \mathcal{M}[\Phi](1) + \mathcal{M}[\Phi](0) \\
        W_{\text{arith}}(\Phi) \;&=\; \sum_{p \in \mathbb{P}} \sum_{n=1}^{\infty} \frac{\ln(p)}{p^{n/2}} \left( \Phi(p^n) + \Phi\left(\frac{1}{p^n}\right) \right) \\
        W_{\text{arch}}(\Phi) \;&=\; \left( \ln(\pi) + \gamma \right) \Phi(1) + \int_{1}^{\infty} \frac{\Phi(x) + \Phi\left(\frac{1}{x}\right) - 2\Phi(1)}{x - \frac{1}{x}} \, \frac{\mathrm{d}x}{x}
    \end{align*}
    with $\mathcal{M}$ denoting the Mellin transform, $\mathbb{P}$ the set of prime numbers and $\gamma$ the Euler-Mascheroni constant.
\end{DefNefEn}

\begin{propositionEn}[label=pr:CriterePositiviteWeil]{Weil positivity criterion}
    Let $\mathcal{Z}$ be the set of non-trivial zeros of the Riemann Zeta function.\par
    According to André Weil's theorem (see\cite{weil1952_formules_en}), the Hermitian form $\langle \cdot, \cdot \rangle_{\mathcal{HP}}$ has the following properties: \par
    For any function $f \in \mathcal{W}$, the evaluation of this form on the diagonal is expressed analytically as a sum over the set $\mathcal{Z}$:
    \begin{equation*}
        \langle f, f \rangle_{\mathcal{HP}} \;=\; \sum_{\rho \in \mathcal{Z}} \mathcal{M}[f](\rho) \, \overline{\mathcal{M}[f](1 - \overline{\rho})}
    \end{equation*}
    where $\mathcal{M}[f]$ denotes the Mellin transform of the test function $f$. \par
    Moreover, the Riemann Hypothesis is satisfied if and only if, for any non-zero function $f \in \mathcal{W}$, we have:
    \begin{equation*}
        \langle f, f \rangle_{\mathcal{HP}} \;\geq\; 0
    \end{equation*}
\end{propositionEn}

\begin{proof}
    \begin{part_proof}
        Let us first prove the equality given in the proposition. \par
        Let $f \in \mathcal{W}$. \par
        Let $\Phi = f * f^*$. By definition of the Weil Hermitian form, we have:
        \begin{equation*}
            \langle f, f \rangle_{\mathcal{HP}} \;=\; W_{\text{pol}}(\Phi) - W_{\text{arith}}(\Phi) - W_{\text{arch}}(\Phi)
        \end{equation*}
        According to André Weil's explicit formulas, the sum of the polar, arithmetic, and Archimedean contributions evaluated on a test function is equal to the sum of its Mellin transform evaluated over the set $\mathcal{Z}$. 
        We deduce that:
        \begin{equation*}
            \langle f, f \rangle_{\mathcal{HP}} \;=\; \sum_{\rho \in \mathcal{Z}} \mathcal{M}[\Phi](\rho)
        \end{equation*}
        By the property of the multiplicative convolution product associated with the Haar measure $\frac{\mathrm{d}x}{x}$, the Mellin transform turns the convolution into a usual product. For any complex number $s$, we have:
        \begin{equation*}
            \mathcal{M}[\Phi](s) \;=\; \mathcal{M}[f](s) \mathcal{M}[f^*](s)
        \end{equation*}
        Let us express the Mellin transform of the involution $f^*$. We have:
        \begin{equation*}
            \mathcal{M}[f^*](s) \;=\; \int_{0}^{\infty} f^*(x) x^{s-1} \, \mathrm{d}x \;=\; \int_{0}^{\infty} \frac{1}{x} \overline{f\left(\frac{1}{x}\right)} x^{s-1} \, \mathrm{d}x
        \end{equation*}
        We perform the change of variable $y = \frac{1}{x}$. We deduce that $\mathrm{d}x = - \frac{1}{y^2} \mathrm{d}y$. Since the integration bounds are reversed, we have:
        \begin{equation*}
            \mathcal{M}[f^*](s) \;=\; \int_{0}^{\infty} y \overline{f(y)} y^{1-s} \frac{1}{y^2} \, \mathrm{d}y \;=\; \int_{0}^{\infty} \overline{f(y)} y^{-s} \, \mathrm{d}y
        \end{equation*}
        Since $y^{-s} = y^{(1 - s) - 1}$, we rewrite the integral in the form:
        \begin{equation*}
            \mathcal{M}[f^*](s) \;=\; \overline{ \int_{0}^{\infty} f(y) y^{(1 - \overline{s}) - 1} \, \mathrm{d}y } \;=\; \overline{\mathcal{M}[f](1 - \overline{s})}
        \end{equation*}
        By evaluating this expression at a non-trivial zero $s = \rho \in \mathcal{Z}$, we obtain:
        \begin{equation*}
            \mathcal{M}[\Phi](\rho) \;=\; \mathcal{M}[f](\rho) \, \overline{\mathcal{M}[f](1 - \overline{\rho})}
        \end{equation*}
        By substituting this identity back into the sum originating from the explicit formulas, we deduce that:
        \begin{equation*}
            \langle f, f \rangle_{\mathcal{HP}} \;=\; \sum_{\rho \in \mathcal{Z}} \mathcal{M}[f](\rho) \, \overline{\mathcal{M}[f](1 - \overline{\rho})}
        \end{equation*}
    \end{part_proof}

    \begin{part_proof}
        Let us now prove the equivalence between the positivity of the Weil Hermitian form and the Riemann hypothesis, by first clarifying the connection with Fourier analysis. \medskip\par
        Assume that the Riemann Hypothesis is satisfied. \par
        In order to make the connection between the Mellin transform and the Fourier transform, we perform a change of variable allowing us to pass from the multiplicative group $\left(\mathbb{R}_+^*, \times\right)$ to the additive group $\left(\mathbb{R}, +\right)$. \par
        Let $x \in \mathbb{R}_+^*$. Let $x = e^u$, with $u \in \mathbb{R}$. \par
        The multiplicative Haar measure then becomes $\frac{\mathrm{d}x}{x} = \mathrm{d}u$. \par
        Let $f \in \mathcal{W}$. Let $F$ be the function defined for all $u \in \mathbb{R}$ by:
        \begin{equation*}
            F(u) \;=\; f\left(e^u\right) e^{\frac{u}{2}}
        \end{equation*}
        Let $\rho \in \mathcal{Z}$. By hypothesis, the real part of $\rho$ is equal to $\frac{1}{2}$. \par
        Let $\rho = \frac{1}{2} + i\gamma$, with $\gamma \in \mathbb{R}$. \par
        Let us express the Mellin transform of $f$ evaluated at $\rho$. By definition, we have:
        \begin{equation*}
            \mathcal{M}[f](\rho) \;=\; \int_{0}^{\infty} f(x) x^{\rho} \, \frac{\mathrm{d}x}{x}
        \end{equation*}
        By applying the change of variable $x = e^u$, we obtain:
        \begin{equation*}
            \mathcal{M}[f](\rho) \;=\; 
            \int_{\mathbb{R}} f\left(e^u\right) e^{\left(\frac{1}{2} + i\gamma\right)u} \, \mathrm{d}u \;=\; 
            \int_{\mathbb{R}} \left( f\left(e^u\right) e^{\frac{u}{2}} \right) e^{i\gamma u} \, \mathrm{d}u
        \end{equation*}
        Since $F(u) = f(e^u) e^{\frac{u}{2}}$, we deduce that:
        \begin{equation*}
            \mathcal{M}[f](\rho) \;=\; \int_{\mathbb{R}} F(u) e^{i\gamma u} \, \mathrm{d}u \;=\; \hat{F}(\gamma)
        \end{equation*}
        where $\hat{F}$ denotes the usual Fourier transform of the function $F$. \par
        Let us now evaluate the conjugate term appearing in the sum of the explicit formula. Since $1 - \overline{\rho} = 1 - \left( \frac{1}{2} - i\gamma \right) = \frac{1}{2} + i\gamma = \rho$, we deduce that:
        \begin{equation*}
            \mathcal{M}[f](1 - \overline{\rho}) \;=\; \mathcal{M}[f](\rho) \;=\; \hat{F}(\gamma)
        \end{equation*}
        Consequently, the product evaluated at $\rho$ is written:
        \begin{equation*}
            \mathcal{M}[f](\rho) \, \overline{\mathcal{M}[f](1 - \overline{\rho})} \;=\; \hat{F}(\gamma) \, \overline{\hat{F}(\gamma)} \;=\; \left| \hat{F}(\gamma) \right|^2
        \end{equation*}
        Since the squared modulus is a positive or zero quantity, we deduce by summation over the set $\mathcal{Z}$ of zeros that:
        \begin{equation*}
            \langle f, f \rangle_{\mathcal{HP}} \;=\; \sum_{\rho \in \mathcal{Z}} \left| \hat{F}(\gamma_\rho) \right|^2 \;\geq\; 0
        \end{equation*}
        The Weil Hermitian form is therefore indeed positive on the space $\mathcal{W}$.
    \end{part_proof}

    \begin{part_proof}
        Conversely, let us suppose that the Riemann Hypothesis is not satisfied. \par
        There exists then a non-trivial zero $\rho_0 \in \mathcal{Z}$ such that, if we set $\rho_0 = \beta_0 + i\gamma_0$, we have $\beta_0 \neq \frac{1}{2}$. \par
        According to the functional symmetry properties of the zeta function, the point $\rho_1 = 1 - \overline{\rho_0} = 1 - \beta_0 + i\gamma_0$ is also a non-trivial zero. \par
        Since $\beta_0 \neq \frac{1}{2}$, we deduce that $\rho_0 \neq \rho_1$. \medskip\par

        According to the works of André Weil (see\cite{weil1952_formules_en}), the space $\mathcal{W}$ possesses the necessary analytic localization properties. \par
        Let $\varepsilon$ be a strictly positive real number such that $\varepsilon < 2$. \par
        It is possible to construct a test function $f \in \mathcal{W}$ such that its Mellin transform satisfies $\mathcal{M}[f](\rho_0) = 1$ and $\mathcal{M}[f](\rho_1) = -1$, while guaranteeing a sufficiently rapid decay so that the sum over the rest of the zeros is bounded in modulus by $\varepsilon$. \par
        Let $\mathcal{Z}' = \mathcal{Z} \setminus \{\rho_0, \rho_1\}$. \par
        We require the function $f$ to satisfy the following inequality:
        \begin{equation*}
            \left| \sum_{\rho \in \mathcal{Z}'} \mathcal{M}[f](\rho) \, \overline{\mathcal{M}[f](1 - \overline{\rho})} \right| \;\leq\; \varepsilon
        \end{equation*}
        We evaluate the sum of the Hermitian form restricted to the two zeros $\rho_0$ and $\rho_1$:
        \begin{align*}
            &\mathcal{M}[f](\rho_0) \, \overline{\mathcal{M}[f](1 - \overline{\rho_0})} \;+\; \mathcal{M}[f](\rho_1) \, \overline{\mathcal{M}[f](1 - \overline{\rho_1})}\\
            &=\; \mathcal{M}[f](\rho_0) \, \overline{\mathcal{M}[f](\rho_1)} \;+\; \mathcal{M}[f](\rho_1) \, \overline{\mathcal{M}[f](\rho_0)}
        \end{align*}
        By replacing with the chosen values, we obtain:
        \begin{equation*}
            1 \cdot \overline{(-1)} \;+\; (-1) \cdot \overline{1} \;=\; -1 \;-\; 1 \;=\; -2
        \end{equation*}
        We then decompose the total Hermitian form over the set $\mathcal{Z}$:
        \begin{equation*}
            \langle f, f \rangle_{\mathcal{HP}} \;=\; -2 \;+\; \sum_{\rho \in \mathcal{Z}'} \mathcal{M}[f](\rho) \, \overline{\mathcal{M}[f](1 - \overline{\rho})}
        \end{equation*}
        Since we imposed that the modulus of the sum over $\mathcal{Z}'$ is less than or equal to $\varepsilon$, we deduce that:
        \begin{equation*}
            \langle f, f \rangle_{\mathcal{HP}} \;\leq\; -2 \;+\; \varepsilon
        \end{equation*}
        Since $\varepsilon < 2$, we deduce that $\langle f, f \rangle_{\mathcal{HP}} < 0$. \par
        The form is therefore not positive. \par
        By contraposition, the positivity of the form $\langle \cdot, \cdot \rangle_{\mathcal{HP}}$ for any function $f \in \mathcal{W}$ implies that all non-trivial zeros satisfy $\beta = \frac{1}{2}$, which proves the Riemann Hypothesis under the assumptions of the proposition.
    \end{part_proof}
\end{proof}


\color{red}
\section{Action of the Weil distribution}
\color{black}

\color{red}
\subsection{von Mangoldt arithmetic function, Möbius function}
\color{black}

This section is dedicated to establishing a specific formulation for 
the Weil Hermitian form, and to detailing the various elements necessary 
for understanding the desired lemma.


\begin{DefNefEn}[label=df:VonMangoldt]{von Mangoldt arithmetic function}
    Following the work of Hans von Mangoldt 
    (see Gérald Tenenbaum~\cite{tenenbaum2015_intro_en}), translating the poles of the logarithmic derivative of the Riemann zeta function requires the introduction of a function supported solely by prime powers. \par
    Let $n \in \mathbb{N}^*$. \par
    We call the \textbf{von Mangoldt function}, denoted by $\Lambda(n)$, the mapping defined by:
    \begin{equation*}
        \Lambda(n) \;=\; 
        \begin{cases}
            \ln(p) & \text{if } n = p^k \text{ with } p \text{ prime and } k \in \mathbb{N}^* \\
            0 & \text{otherwise}
        \end{cases}
    \end{equation*}
\end{DefNefEn}

\begin{DefNefEn}[label=df:MobiusInversion]{Möbius function and inversion formula}
    According to the foundations of analytic number theory 
    (see Gérald Tenenbaum \cite{tenenbaum2015_intro_en}, Chapter I), we introduce a central arithmetic function for the study of Dirichlet series and arithmetic combinatorics. \par
    Let $n \in \mathbb{N}^*$. \par
    We call the \textbf{Möbius function}, denoted by $\mu(n)$, the mapping defined by:
    \begin{equation*}
        \mu(n) \;=\; 
        \begin{cases}
            1 & \text{if } n = 1 \\
            (-1)^r & \text{if } n \text{ is the product of } r \text{ distinct prime numbers} \\
            0 & \text{if } n \text{ has a strict square factor}
        \end{cases}
    \end{equation*}
    We call the \textbf{Möbius inversion formula} the following result: for all arithmetic functions $f$ and $g$, we have the strict equivalence:
    \begin{equation*}
        g(n) \;=\; \sum_{d \mid n} f(d) \quad \iff \quad f(n) \;=\; \sum_{d \mid n} \mu(d) \, g\left(\frac{n}{d}\right)
    \end{equation*}
\end{DefNefEn}


\begin{propositionEn}[label=pr:LienMangoldtLog]{Arithmetic link between the von Mangoldt function and the logarithm}
    Let $n \in \mathbb{N}^*$. \par
    The von Mangoldt function $\Lambda$ satisfies the following summation identity:
    \begin{equation*}
        \sum_{d \mid n} \Lambda(d) \;=\; \ln(n)
    \end{equation*}
    Consequently, by applying the Möbius inversion formula, we obtain the explicit expression:
    \begin{equation*}
        \Lambda(n) \;=\; \sum_{d \mid n} \mu(d) \ln\left(\frac{n}{d}\right) \;=\; -\sum_{d \mid n} \mu(d) \ln(d)
    \end{equation*}
\end{propositionEn}

\begin{proof}
    \begin{part_proof}
        Let $n \in \mathbb{N}^*$. \par
        If $n=1$, we have on the one hand $\sum_{d \mid 1} \Lambda(d) = \Lambda(1) = 0$, and on the other hand $\ln(1) = 0$. The equality is therefore verified. \par
        Assume $n > 1$. \par
        Let the prime factorization of the integer $n$ be:
        \begin{equation*}
            n \;=\; \prod_{i=1}^{k} p_i^{\alpha_i}
        \end{equation*}
        where the $p_i$ are distinct prime numbers and the $\alpha_i$ are strictly positive integers. \medskip\par
        Since we evaluate the sum $\sum_{d \mid n} \Lambda(d)$ and, by definition, the function $\Lambda$ is zero except on strict prime powers, we deduce that the only divisors $d$ of $n$ providing a non-zero contribution to the sum are of the form $d = p_i^j$ with $1 \leq j \leq \alpha_i$. \par
        We then have:
        \begin{equation*}
            \sum_{d \mid n} \Lambda(d) \;=\; \sum_{i=1}^{k} \sum_{j=1}^{\alpha_i} \Lambda\left(p_i^j\right)
        \end{equation*}
        Since $\Lambda(p_i^j) = \ln(p_i)$ for all $j \geq 1$, we deduce that:
        \begin{equation*}
            \sum_{d \mid n} \Lambda(d) \;=\; \sum_{i=1}^{k} \alpha_i \ln(p_i) \;=\; \ln\left( \prod_{i=1}^{k} p_i^{\alpha_i} \right) \;=\; \ln(n)
        \end{equation*}
        We have thus proved the first equality. \par
        Next, let $f = \Lambda$ and $g = \ln$. Since $g(n) = \sum_{d \mid n} f(d)$, we deduce by the Möbius inversion formula (Definition~\ref{df:MobiusInversion}) that:
        \begin{equation*}
            \Lambda(n) \;=\; \sum_{d \mid n} \mu(d) \ln\left(\frac{n}{d}\right)
        \end{equation*}
        Since $\ln\left(\frac{n}{d}\right) = \ln(n) - \ln(d)$, and knowing that an elementary property of the Möbius function gives $\sum_{d \mid n} \mu(d) = 0$ for all $n > 1$, we deduce that:
        \begin{equation*}
            \sum_{d \mid n} \mu(d) \ln(n) \;=\; \ln(n) \sum_{d \mid n} \mu(d) \;=\; 0
        \end{equation*}
        We finally deduce the equality:
        \begin{equation*}
            \Lambda(n) \;=\; -\sum_{d \mid n} \mu(d) \ln(d)
        \end{equation*}
        which completes the proof.
    \end{part_proof}
\end{proof}


\color{red}
\subsection{Archimedean kernel}
\color{black}

\begin{DefNefEn}[label=df:FonctionGammaEuler]{Euler's Gamma function}
    Following the work of Leonhard Euler on the factorial extension, 
    we introduce the following parametric integral: \par
    Let $z \in \mathbb{C}$ such that $\Re(z) > 0$. \par
    We call \textbf{Euler's Gamma function}, denoted by $\Gamma(z)$, 
    the function defined by the absolutely convergent integral:
    \begin{equation*}
        \Gamma(z) \;=\; \int_{0}^{+\infty} t^{z-1} e^{-t} \,\mathrm{d}t
    \end{equation*}
\end{DefNefEn}

\begin{DefNefEn}[label=df:NoyauArchimedien]{Archimedean kernel and local factor at infinity}
    Following the theory developed by Bernhard Riemann and later generalized by André Weil 
    (see~\cite{weil1952_formules_en}), 
    the functional equation of the zeta function requires completing it with a local factor associated with the Archimedean place, that is to say linked to the completion of the field $\mathbb{Q}$ for the usual absolute value giving $\mathbb{R}$. \par
    Let $s \in \mathbb{C}$ such that $\Re(s) > 0$. \par
    We define the Gamma factor, or Archimedean local factor, associated with the Riemann zeta function:
    \begin{equation*}
        \Gamma_{\mathbb{R}}(s) \;=\; \pi^{-\frac{s}{2}} \Gamma\left(\frac{s}{2}\right)
    \end{equation*}
    Let $x$ be a strictly positive real number and $y$ a strictly positive real number. \par
    We equip the multiplicative group $\mathbb{R}_+^*$ with its invariant Haar measure, given by $\frac{\mathrm{d}t}{t}$. \par
    We call the \textbf{Archimedean kernel}, denoted by $\mathcal{K}_{\infty}(x,y)$, the distribution on the multiplicative group $\mathbb{R}_+^*$ defined by the multiplicative convolution action $\mathcal{K}_{\infty}(x,y) = W_{\infty}\left(\frac{x}{y}\right)$, where $W_{\infty}$ is the distribution on $\mathbb{R}_+^*$ whose Mellin transform corresponds to the opposite of the logarithmic derivative of the Gamma factor:
    \begin{equation*}
        \mathcal{M}[W_{\infty}](s) \;=\; -\frac{\Gamma_{\mathbb{R}}'(s)}{\Gamma_{\mathbb{R}}(s)}
    \end{equation*}
\end{DefNefEn}

\begin{DefNefEn}[label=df:FonctionDigamma]{Digamma function and logarithmic derivative}
    According to the fundamental properties of complex analysis, the composition of a meromorphic function with the complex logarithm requires fixing a branch. To circumvent this topological obstruction, we define the logarithmic derivative algebraically. \par
    Let $U$ be a connected open set of $\mathbb{C}$. \par
    Let $f$ be a meromorphic function on $U$, not identically zero. \par
    We call the \textbf{logarithmic derivative} of $f$, the meromorphic function on $U$ defined by the quotient $\frac{f'}{f}$. \par
    Let $z \in \mathbb{C}$ such that $z$ is not a negative integer or zero. \par
    We call the \textbf{Digamma function}, denoted by $\psi(z)$, the logarithmic derivative of Euler's Gamma function:
    \begin{equation*}
        \psi(z) \;=\; \frac{\Gamma'(z)}{\Gamma(z)}
    \end{equation*}
\end{DefNefEn}

\begin{propositionEn}[label=pr:CalculDeriveeLogArchimedienne]{Evaluation of the Archimedean meromorphic symbol}
    Let $s \in \mathbb{C}$ such that $\Re(s) > 0$. \par
    The opposite of the logarithmic derivative of the Archimedean local factor $\Gamma_{\mathbb{R}}(s) = \pi^{-\frac{s}{2}} \Gamma\left(\frac{s}{2}\right)$ admits the explicit expression:
    \begin{equation*}
        -\frac{\Gamma_{\mathbb{R}}'(s)}{\Gamma_{\mathbb{R}}(s)} \;=\; \frac{1}{2}\ln(\pi) - \frac{1}{2}\psi\left(\frac{s}{2}\right)
    \end{equation*}
\end{propositionEn}

\begin{proof}
    \begin{part_proof}
        Let $s \in \mathbb{C}$ such that $\Re(s) > 0$. \par
        Let us set the functions $u(s) = \pi^{-\frac{s}{2}}$ and $v(s) = \Gamma\left(\frac{s}{2}\right)$. \par
        Since we have the equality $\pi^{-\frac{s}{2}} = e^{-\frac{s}{2} \ln(\pi)}$, we deduce by standard differentiation of exponential functions that:
        \begin{equation*}
            u'(s) \;=\; -\frac{1}{2}\ln(\pi) \pi^{-\frac{s}{2}}
        \end{equation*}
        On the other hand, by the chain rule, we have:
        \begin{equation*}
            v'(s) \;=\; \frac{1}{2} \Gamma'\left(\frac{s}{2}\right)
        \end{equation*}
        Since $\Gamma_{\mathbb{R}}(s) = u(s)v(s)$, we deduce by the Leibniz rule that:
        \begin{equation*}
            \Gamma_{\mathbb{R}}'(s) \;=\; u'(s)v(s) + u(s)v'(s) \;=\; -\frac{1}{2}\ln(\pi) \pi^{-\frac{s}{2}} \Gamma\left(\frac{s}{2}\right) + \pi^{-\frac{s}{2}} \frac{1}{2} \Gamma'\left(\frac{s}{2}\right)
        \end{equation*}
        By dividing this relation by the initial expression of $\Gamma_{\mathbb{R}}(s)$, we obtain the quotient:
        \begin{equation*}
            \frac{\Gamma_{\mathbb{R}}'(s)}{\Gamma_{\mathbb{R}}(s)} \;=\; \frac{-\frac{1}{2}\ln(\pi) \pi^{-\frac{s}{2}} \Gamma\left(\frac{s}{2}\right)}{\pi^{-\frac{s}{2}} \Gamma\left(\frac{s}{2}\right)} + \frac{\frac{1}{2} \pi^{-\frac{s}{2}} \Gamma'\left(\frac{s}{2}\right)}{\pi^{-\frac{s}{2}} \Gamma\left(\frac{s}{2}\right)}
        \end{equation*}
        By simplifying the strictly non-zero terms, the equation becomes:
        \begin{equation*}
            \frac{\Gamma_{\mathbb{R}}'(s)}{\Gamma_{\mathbb{R}}(s)} \;=\; -\frac{1}{2}\ln(\pi) + \frac{1}{2} \frac{\Gamma'\left(\frac{s}{2}\right)}{\Gamma\left(\frac{s}{2}\right)}
        \end{equation*}
        Since we defined $\psi(z) = \frac{\Gamma'(z)}{\Gamma(z)}$ (Definition~\ref{df:FonctionDigamma}), we deduce by evaluating at the point $z = \frac{s}{2}$ that:
        \begin{equation*}
            \frac{\Gamma_{\mathbb{R}}'(s)}{\Gamma_{\mathbb{R}}(s)} \;=\; -\frac{1}{2}\ln(\pi) + \frac{1}{2}\psi\left(\frac{s}{2}\right)
        \end{equation*}
        By taking the opposite of this last expression, we finally deduce that:
        \begin{equation*}
            -\frac{\Gamma_{\mathbb{R}}'(s)}{\Gamma_{\mathbb{R}}(s)} \;=\; \frac{1}{2}\ln(\pi) - \frac{1}{2}\psi\left(\frac{s}{2}\right)
        \end{equation*}
        This result justifies the expansion of the analytic kernel in the Mellin transform.
    \end{part_proof}
\end{proof}

\begin{corollaireEn}[label=cr:EvaluationMellinNoyauArchimedien]{Explicit evaluation of the Mellin transform of the Archimedean kernel}
    Let $s \in \mathbb{C}$ such that $\Re(s) > 0$. \par
    The Mellin transform of the distribution $W_{\infty}$ is expressed using the constant $\pi$ and the Digamma function $\psi$ by the equality:
    \begin{equation*}
        \mathcal{M}[W_{\infty}](s) \;=\; \frac{1}{2}\ln(\pi) - \frac{1}{2}\psi\left(\frac{s}{2}\right)
    \end{equation*}
\end{corollaireEn}

\begin{proof}
    \begin{part_proof}
        By definition of the Archimedean kernel (Definition~\ref{df:NoyauArchimedien}), we have the equality:
        \begin{equation*}
            \mathcal{M}[W_{\infty}](s) \;=\; -\frac{\Gamma_{\mathbb{R}}'(s)}{\Gamma_{\mathbb{R}}(s)}
        \end{equation*}
        Since we established in Proposition~\ref{pr:CalculDeriveeLogArchimedienne} the evaluation of this quotient, we deduce by directly substituting the expression that:
        \begin{equation*}
            \mathcal{M}[W_{\infty}](s) \;=\; \frac{1}{2}\ln(\pi) - \frac{1}{2}\psi\left(\frac{s}{2}\right)
        \end{equation*}
        Since we defined the kernel $\mathcal{K}_{\infty}$ on the group $(\mathbb{R}_+^*, \times)$ without analytic ambiguity, we deduce that its action by integration against a test function of the space $\mathcal{W}$ with respect to the invariant Haar measure exactly yields the contribution of the place at infinity in the framework of Weil's explicit formula.
    \end{part_proof}
\end{proof}


\color{red}
\subsection{Weil distribution action lemma}
\color{black}

\begin{DefNefEn}[label=df:DistributionDiracMultiplicative]{Dirac distribution on the multiplicative group}
    Following the foundations of functional analysis laid by Laurent Schwartz 
    (see~\cite{schwartz1966_distributions_en}), we define the Dirac distribution adapted to the multiplicative group $\mathbb{R}_+^*$. \par
    Let $a$ be a strictly positive real number. \par
    We call the \textbf{Dirac distribution centered at $a$}, denoted by $\delta_a$ or abusively by the kernel $x \mapsto \delta(x-a)$ in the integral notation, the distribution acting on the test function space $\mathcal{W}$ such that, for any function $f \in \mathcal{W}$, we have:
    \begin{equation*}
        \langle \delta_a, f \rangle \;=\; \int_{0}^{\infty} f(x) \, \delta(x-a) \, \frac{\mathrm{d}x}{x} \;:=\; f(a)
    \end{equation*}
    By extension to the tensor product on the space $\mathbb{R}_+^* \times \mathbb{R}_+^*$, we define the \textbf{diagonal Dirac kernel}, denoted by $\delta(x-y)$, acting on two test functions $f \in \mathcal{W}$ and $g \in \mathcal{W}$ by the duality relation:
    \begin{equation*}
        \int_{0}^{\infty} \int_{0}^{\infty} f(x) \overline{g(y)} \, \delta(x-y) \, \frac{\mathrm{d}x}{x} \frac{\mathrm{d}y}{y} \;:=\; \int_{0}^{\infty} f(x) \overline{g(x)} \, \frac{\mathrm{d}x}{x}
    \end{equation*}
    This formalism allows defining the pointwise components of bilinear distributions.
\end{DefNefEn}

\begin{lemmaEn}[label=lm:EquivalenceFormeWeil]{Action of the Weil distribution and spectral equivalence}
    Let $f$ and $g$ be two functions in the Weil space $\mathcal{W}$. \par
    Let us define the Weil distribution $\mathcal{K}_{\Gamma}$ on the product $\mathbb{R}_+^* \times \mathbb{R}_+^*$ by the sum of its components (identity, arithmetic, and Archimedean):
    \begin{equation*}
        \mathcal{K}_{\Gamma}(x,y) \;=\; \delta(x-y) \;-\; \sum_{n=1}^{\infty} \frac{\Lambda(n)}{\sqrt{n}} \left[ \delta\left(x - ny\right) + \delta\left(y - nx\right) \right] \;+\; \mathcal{K}_{\infty}(x,y)
    \end{equation*}
    where $\Lambda$ is the von Mangoldt function (Definition~\ref{df:VonMangoldt}), $\mathcal{K}_{\infty}$ is the Archimedean kernel (Definition~\ref{df:NoyauArchimedien}), and $\delta$ is the Dirac distribution (Definition~\ref{df:DistributionDiracMultiplicative}). \par
    We define the action of this distribution on the tensor product $f \otimes \overline{g}$ with respect to the invariant multiplicative Haar measure $\frac{\mathrm{d}x}{x}$:
    \begin{equation*}
        \langle \mathcal{K}_{\Gamma}, f \otimes \overline{g} \rangle \;:=\; \int_{0}^{\infty} \int_{0}^{\infty} f(x) \overline{g(y)} \, \mathcal{K}_{\Gamma}(x,y) \, \frac{\mathrm{d}x}{x} \frac{\mathrm{d}y}{y}
    \end{equation*}
    According to the work of André Weil (see~\cite{weil1952_formules_en}), this double integral converges absolutely and satisfies the following spectral equality on the set $\mathcal{Z}$ of non-trivial zeros of the Riemann zeta function:
    \begin{equation*}
        \langle \mathcal{K}_{\Gamma}, f \otimes \overline{g} \rangle \;=\; \sum_{\rho \in \mathcal{Z}} \mathcal{M}[f](\rho) \, \overline{\mathcal{M}[g](1 - \overline{\rho})}
    \end{equation*}
    Consequently, the action of the distribution exactly coincides with the associated Weil sesquilinear form $\langle f, g \rangle_{\mathcal{HP}}$, described in {Proposition}~\ref{pr:CriterePositiviteWeil}, and defined in~\ref{df:FormeHermitienneWeil}.
\end{lemmaEn}

\begin{proof}
    \begin{part_proof}
        Let $f \in \mathcal{W}$ and $g \in \mathcal{W}$. \par
        Let us define the function $h$ on the multiplicative group $\mathbb{R}_+^*$ by the multiplicative convolution:
        \begin{equation*}
            h(u) \;=\; \int_{0}^{\infty} f(uy) \, \overline{g(y)} \, \frac{\mathrm{d}y}{y}
        \end{equation*}
        Since the space $\mathcal{W}$ is stable under multiplicative convolution and conjugation, we deduce that $h \in \mathcal{W}$. \par
        According to the properties of the Mellin transform on convolution products, the transform of the function $h$ evaluated at a complex point $s$ satisfies the equality:
        \begin{equation*}
            \mathcal{M}[h](s) \;=\; \mathcal{M}[f](s) \, \overline{\mathcal{M}[g](1 - \overline{s})}
        \end{equation*}
        By expanding the double integral defining the action of the distribution:
        \begin{equation*}
            \langle \mathcal{K}_{\Gamma}, f \otimes \overline{g} \rangle \;=\; \int_{0}^{\infty} \int_{0}^{\infty} f(x) \overline{g(y)} \, \mathcal{K}_{\Gamma}(x,y) \, \frac{\mathrm{d}x}{x} \frac{\mathrm{d}y}{y}
        \end{equation*}
        We apply the change of variable $x = uy$. \par
        Since the Haar measure $\frac{\mathrm{d}x}{x}$ is invariant under dilation, we have $\frac{\mathrm{d}x}{x} = \frac{\mathrm{d}u}{u}$. The action of the bilinear distribution $\mathcal{K}_{\Gamma}$ on the function $(u, y) \mapsto f(uy)\overline{g(y)}$ exactly reduces to the evaluation of the one-dimensional distribution of the explicit formula, denoted by $W$, on the test function $h$. \par
        According to the Weil explicit formula theorem, this geometric 
        evaluation is equal to the sum of the mellin transforme over the set $\mathcal{Z}$:
        \begin{equation*}
            \langle \mathcal{K}_{\Gamma}, f \otimes \overline{g} \rangle \;=\; \langle W, h \rangle_{\text{Explicite}} \;=\; \sum_{\rho \in \mathcal{Z}} \mathcal{M}[h](\rho)
        \end{equation*}
        By substituting the expression for the Mellin transform of $h$ obtained previously, we get the equality:
        \begin{equation*}
            \langle \mathcal{K}_{\Gamma}, f \otimes \overline{g} \rangle \;=\; \sum_{\rho \in \mathcal{Z}} \mathcal{M}[f](\rho) \, \overline{\mathcal{M}[g](1 - \overline{\rho})}
        \end{equation*}
        Since the functions $f$ and $g$ belong to the Weil space $\mathcal{W}$, their Mellin transforms are rapidly decreasing. \par
        We deduce that the series indexed by the zeros $\rho \in \mathcal{Z}$ is absolutely convergent. \par
        The spectral equivalence is thus demonstrated, and we thereby justify the existence of the sesquilinear form $\langle f, g \rangle_{\mathcal{HP}}$ for the Hilbert-Pólya space.
    \end{part_proof}
\end{proof}


\begin{DefNefEn}[label=df:EspaceHilbertPolya]{Hilbert-Pólya space}
    According to the historical conjecture of David Hilbert and George Pólya, 
    whose concrete realization was formalized by Alain Connes (see~\cite{connes1999_trace_en} (Section IV)), 
    we introduce a Hilbert space associated with the distribution of the zeros of the 
    Riemann zeta function. \par
    Let $\mathcal{W}$ be the space of Weil test functions. \par
    Let $\mathcal{N}$ be the vector subspace of $\mathcal{W}$ corresponding to the isotropic kernel of the Weil Hermitian form, defined by:
    \begin{equation*}
        \mathcal{N} \;:=\; \left\{ f \in \mathcal{W} \mid \langle f, f \rangle_{\mathcal{HP}} = 0 \right\}
    \end{equation*}
    Under the Riemann Hypothesis, according to {Proposition}~\ref{pr:CriterePositiviteWeil}, the form $\langle \cdot, \cdot \rangle_{\mathcal{HP}}$ is positive semi-definite. \par
    Since the Weil form is a positive Hermitian form, we deduce by passing to the quotient that it induces a positive definite inner product on the quotient vector space $\mathcal{W} / \mathcal{N}$. \par
    We deduce that the space $\mathcal{W} / \mathcal{N}$, equipped with the associated metric, has the structure of a separated pre-Hilbert space. \par
    We call the \textbf{Hilbert-Pólya space}, denoted by $\mathcal{H}_{HP}$, the topological completion of this quotient space with respect to the induced norm $\|\cdot\|_{\mathcal{HP}}$:
    \begin{equation*}
        \mathcal{H}_{HP} \;:=\; \overline{\left( \mathcal{W} / \mathcal{N} \right)}^{\|\cdot\|_{\mathcal{HP}}}
    \end{equation*}
\end{DefNefEn}

\begin{lemmaEn}[label=lm:ConstructionMetriqueHP]{Hilbert-Pólya structure and Hamiltonian operator}
    Let $\mathcal{Z}$ be the set of non-trivial zeros of the Riemann Zeta function.\par
    According to Lemma~\ref{lm:EquivalenceFormeWeil}, 
    the action of the Weil bilinear distribution on the space $\mathcal{W}$ 
    coincides with the spectral sesquilinear form, which we 
    will denote by $\langle \cdot, \cdot \rangle_{\mathcal{HP}}$. \par
    According to {Proposition}~\ref{pr:CriterePositiviteWeil} (Weil positivity 
    criterion), this Hermitian form satisfies $\langle f, f \rangle_{\mathcal{HP}} 
    \geq 0$ for any function $f \in \mathcal{W}$ if and only if the Riemann 
    Hypothesis is true. \par
    Under this condition, let $\mathcal{N}$ be the kernel associated with this form 
    (isotropic subspace). \par
    Since the form $\langle \cdot, \cdot \rangle_{\mathcal{HP}}$ induces a 
    strictly positive definite inner product on the quotient space 
    $\mathcal{W} / \mathcal{N}$, we deduce that the topological completion of this 
    quotient space generates a separable Hilbert space, which we will denote by 
    $\mathcal{H}_{HP}$. \par
    On this Hilbert space $\mathcal{H}_{HP}$, there exists an unbounded 
    spectral Hamiltonian operator, which we will denote by $\mathbf{H}$, whose 
    eigenvalues correspond exactly to the imaginary parts of the non-trivial 
    zeros of the zeta function. \par
    This operator $\mathbf{H}$ is essentially self-adjoint.
\end{lemmaEn}

\begin{proof}
    \begin{part_proof}
        \textbf{Step 1: Characterization of the kernel and definition of the spectral evaluation.} \par
        Assume that the Riemann Hypothesis is true. \par
        According to {Proposition}~\ref{pr:CriterePositiviteWeil}, for any function $f \in \mathcal{W}$, the Weil Hermitian form is non-negative, and we have the equality:
        \begin{equation*}
            \langle f, f \rangle_{\mathcal{HP}} \;=\; \sum_{\rho \in \mathcal{Z}} \left| \mathcal{M}[f](\rho) \right|^2
        \end{equation*}
        According to Definition~\ref{df:EspaceHilbertPolya}, the Hilbert-Pólya space $\mathcal{H}_{HP}$ is constructed by completion of the quotient $\mathcal{W} / \mathcal{N}$, where $\mathcal{N}$ is the isotropic kernel of this Hermitian form. \par
        Since a sum of non-negative real numbers is zero if and only if each term is zero, we deduce from the previous equality the explicit characterization of the kernel:
        \begin{equation*}
            \mathcal{N} \;=\; \left\{ f \in \mathcal{W} \mid \forall \rho \in \mathcal{Z}, \; \mathcal{M}[f](\rho) = 0 \right\}
        \end{equation*}
        Let us define the spectral evaluation map $\mathcal{E}$ on the quotient space $\mathcal{W} / \mathcal{N}$. \par
        Let $[f]$ be a class in the quotient space $\mathcal{W} / \mathcal{N}$ and $f_1, f_2 \in \mathcal{W}$ two representatives of this class. \par
        Since $f_1 - f_2 \in \mathcal{N}$, we deduce by the characterization of the kernel that, for all $\rho \in \mathcal{Z}$, we have $\mathcal{M}[f_1 - f_2](\rho) = 0$. \par
        By linearity of the Mellin transform, we deduce that $\mathcal{M}[f_1](\rho) = \mathcal{M}[f_2](\rho)$. \par
        We deduce that the sequence of spectral evaluations does not depend on the chosen representative. \par
        We thus define the map $\mathcal{E}$ from $\mathcal{W} / \mathcal{N}$ taking values in the space of sequences, defined by:
        \begin{equation*}
            \mathcal{E}([f]) \;:=\; \left( \mathcal{M}[f](\rho) \right)_{\rho \in \mathcal{Z}}
        \end{equation*}
    \end{part_proof}

    \begin{part_proof}
        \textbf{Step 2: Isometric isomorphism with the space of square-summable sequences.} \par
        Let $[f]$ be a class in the quotient space $\mathcal{W} / \mathcal{N}$, represented by the test function $f \in \mathcal{W}$. \par
        Since the function $f$ belongs to the Weil space $\mathcal{W}$, its Mellin transform $\mathcal{M}[f]$ is a rapidly decreasing function on the critical strip. \par
        According to the work of Edward Charles Titchmarsh on the 
        Riemann-von Mangoldt formula (see Gérald Tenenbaum~\cite{tenenbaum2015_intro_en} (Chapter II.3)), the number of zeros $N(T)$ of the zeta function with imaginary parts between $0$ and $T$ grows asymptotically according to the equivalent $\frac{T}{2\pi} \ln\left(\frac{T}{2\pi}\right)$. \par
        Since the rapid decrease of the function $\rho \mapsto |\mathcal{M}[f](\rho)|^2$ decays faster than the inverse of any polynomial, it compensates for the quasi-linear density of the zeros. \par
        We deduce that the series of squared moduli is absolutely convergent on $\mathcal{Z}$. \par
        We deduce that the sequence of its evaluations on $\mathcal{Z}$ is square-summable, and therefore that $\mathcal{E}([f])$ belongs to the space $\ell^2(\mathcal{Z})$. \par
        Furthermore, by definition of the usual norm on $\ell^2(\mathcal{Z})$ and by Proposition~\ref{pr:CriterePositiviteWeil}, we have the equality:
        \begin{equation*}
            \|\mathcal{E}([f])\|_{\ell^2}^2 \;=\; \sum_{\rho \in \mathcal{Z}} \left| \mathcal{M}[f](\rho) \right|^2 \;=\; \langle [f], [f] \rangle_{\mathcal{HP}}
        \end{equation*}
        Since the map $\mathcal{E}$ preserves the metric and respects the vector space structure, we deduce that $\mathcal{E}$ is a linear isometry. \par
        According to the work of Alain Connes on the spectral framework of noncommutative geometry (see~\cite{connes1999_trace_en} (Section IV)), 
        the spectral evaluations of the functions of the Weil space form a dense subspace in $\ell^2(\mathcal{Z})$. \par
        We deduce that the image of the isometry $\mathcal{E}$ is dense in $\ell^2(\mathcal{Z})$. \par
        According to the extension theorem for uniformly continuous maps, since $\mathcal{E}$ is an isometry defined on the pre-Hilbert space $\mathcal{W}/\mathcal{N}$ and taking values in the complete space $\ell^2(\mathcal{Z})$, it uniquely extends to an isometry on the topological completion $\mathcal{H}_{HP}$. \par
        Since its initial image is dense in $\ell^2(\mathcal{Z})$ and the codomain is complete, we deduce that this extension is surjective. \par
        We deduce that the Hilbert-Pólya space $\mathcal{H}_{HP}$ is canonically isometric to the separable Hilbert space $\ell^2(\mathcal{Z})$.
    \end{part_proof}

    \begin{part_proof}
        \textbf{Step 3: Spectral Hamiltonian operator and self-adjointness.} \par
        Let $\mathbf{H}$ be the Hamiltonian operator defined on the dense subspace $\mathcal{W} / \mathcal{N}$ via the isomorphism $\mathcal{E}$ by the relation:
        \begin{equation*}
            \mathcal{E}\left(\mathbf{H} [f]\right)(\rho) \;=\; \gamma_\rho \cdot \mathcal{E}([f])(\rho), \quad \text{where } \rho = \frac{1}{2} + i\gamma_\rho \in \mathcal{Z}
        \end{equation*}
        Let $[f]$ and $[g]$ be two classes in $\mathcal{W} / \mathcal{N}$. \par
        We evaluate the sesquilinear form:
        \begin{equation*}
            \langle \mathbf{H} [f], [g] \rangle_{\mathcal{HP}} \;=\; \sum_{\rho \in \mathcal{Z}} \left( \gamma_\rho \mathcal{M}[f](\rho) \right) \overline{\mathcal{M}[g](\rho)}
        \end{equation*}
        Assume that the Riemann Hypothesis is true. \par
        Under this hypothesis, for all $\rho \in \mathcal{Z}$, we have $\gamma_\rho \in \mathbb{R}$. \par
        Since $\overline{\gamma_\rho} = \gamma_\rho$, we deduce the equality:
        \begin{equation*}
            \langle \mathbf{H} [f], [g] \rangle_{\mathcal{HP}} \;=\; \sum_{\rho \in \mathcal{Z}} \mathcal{M}[f](\rho) \overline{\left( \gamma_\rho \mathcal{M}[g](\rho) \right)} \;=\; \langle [f], \mathbf{H} [g] \rangle_{\mathcal{HP}}
        \end{equation*}
        Since this equality holds for all classes $[f]$ and $[g]$ of the dense domain, we deduce that the operator $\mathbf{H}$ is a symmetric operator. \par
        According to the work of John von Neumann on the spectral theorem for unbounded 
        operators (see\cite{vonneumann_quantum_en} (Chapter II, Section 9)), the multiplication operator by a real sequence, defined on a maximal dense domain in the Hilbert space of square-summable sequences, is essentially self-adjoint. \par
        Since the isometry $\mathcal{E}$ canonically identifies the Hilbert-Pólya space $\mathcal{H}_{HP}$ with a closed subspace of the space $\ell^2(\mathcal{Z})$, we deduce that the operator $\mathbf{H}$ admits a unique self-adjoint extension on $\mathcal{H}_{HP}$. \par
        The complete construction thus demonstrates that the spectral realization of the non-trivial zeros of the Riemann zeta function, by means of a self-adjoint operator, intrinsically relies on the strict positivity of the Weil Hermitian form on the quotient space.
    \end{part_proof}

\end{proof}

\color{red}
\section{Proof of the Riemann hypothesis}
\color{black}

The following definitions are provided here as necessary reminders 
for the theorem that follows:

\begin{DefNefEn}[label=df:OperateurAutoAdjointResolvante]{Self-adjoint operator}
    Let $\mathcal{H}$ be a Hilbert space over the field of complex numbers. \par
    Let $D(A)$ be a dense vector subspace in $\mathcal{H}$ and $A : D(A) \to \mathcal{H}$ be a linear operator, potentially unbounded. \par
    According to the work of Tosio Kato~\cite{kato1995_perturbation_en}, 
    we define the adjoint $A^*$ as follows. \par
    \vspace{0.2cm}
    \textbf{1. Self-adjoint operator:} \par
    Let $D(A^*)$ be the domain of the adjoint operator, defined by the set:
    \begin{equation*}
        D(A^*) \;=\; \left\{ y \in \mathcal{H} : \exists z \in \mathcal{H}, \forall x \in D(A), \langle Ax, y \rangle = \langle x, z \rangle \right\}
    \end{equation*}
    For any $y \in D(A^*)$, by density of $D(A)$, the vector $z$ is unique, and we set $A^*y = z$. \par
    We say that the operator $A$ is \textbf{self-adjoint} if and only if the domains are strictly equal, $D(A) = D(A^*)$, and for any $x \in D(A)$, the equality $Ax = A^*x$ holds. \par
\end{DefNefEn}


\begin{DefNefEn}[label=df:TraceRegularisee]{Spectral regularized trace}
    Let $\mathcal{H}$ be a separable Hilbert space. \par
    Let $\mathcal{L}^1(\mathcal{H})$ be the two-sided ideal of trace-class operators on $\mathcal{H}$, defined as the set of bounded operators for which the usual trace $\mathrm{Tr}$ is absolutely convergent. \par
    According to the work of Alain Connes on geometric trace formulas 
    (see~\cite{connes1999_trace_en}), the study of unbounded spectral operators whose density of states diverges requires an extension of this trace. \par
    Let $\mathcal{A}$ be a subalgebra of operators on $\mathcal{H}$ such that the inclusion $\mathcal{L}^1(\mathcal{H}) \subset \mathcal{A}$ holds. \par
    We call the \textbf{regularized trace}, denoted by $\mathrm{Tr}_{\mathrm{reg}}$, a linear form defined on the algebra $\mathcal{A}$ and satisfying the following compatibility condition: for any operator $T \in \mathcal{L}^1(\mathcal{H})$, we have the equality:
    \begin{equation*}
        \mathrm{Tr}_{\mathrm{reg}}(T) \;=\; \mathrm{Tr}(T)
    \end{equation*}
    For a spectral operator $\mathbf{D}$ with a discrete spectrum but whose number of eigenvalues grows too rapidly, let $\Lambda$ be a strictly positive real number and let $P_\Lambda$ denote the orthogonal spectral projector associated with the interval $[-\Lambda, \Lambda]$. \par
    Let $T \in \mathcal{A}$ be an operator of the form $T = f(\mathbf{D})$. \par
    Since the operator $T$ is not necessarily trace-class on $\mathcal{H}$, we deduce that the quantity $\mathrm{Tr}(T P_\Lambda)$ diverges as $\Lambda$ tends to $+\infty$. \par
    We then define the regularized trace by extracting the finite part of the asymptotic limit:
    \begin{equation*}
        \mathrm{Tr}_{\mathrm{reg}}(T) \;=\; \lim_{\Lambda \to +\infty} \left( \mathrm{Tr}(T P_\Lambda) - \mathcal{C}(T, \Lambda) \right)
    \end{equation*}
    where $\mathcal{C}(T, \Lambda)$ denotes the counter-divergence term (or asymptotic volume), depending continuously on the function $f$, constructed such that the limit exists and defines a geometric invariant on $\mathcal{A}$.
\end{DefNefEn}

\begin{DefNefEn}[label=df:IdealDixmierMacaev]{Macaev Ideal and Dixmier Trace}
    Let $\mathcal{H}$ be a separable Hilbert space. \par
    Let $T$ be a compact operator acting on $\mathcal{H}$, and let 
    $(\mu_n(T))_{n \in \mathbb{N}^*}$ be the sequence of its singular values sorted 
    in decreasing order. \par
    According to the work of Jacques Dixmier \cite{dixmier1966_trace_en}, 
    the Macaev ideal, denoted by $\mathcal{L}^{1,\infty}(\mathcal{H})$, is defined as 
    the set of compact operators $T$ such that:
    \begin{equation*}
        \sup_{N \geq 2} \frac{1}{\ln(N)} \sum_{n=1}^{N} \mu_n(T) \;<\; +\infty
    \end{equation*}
    For any positive operator $T \in \mathcal{L}^{1,\infty}(\mathcal{H})$, the Dixmier 
    trace $\mathrm{Tr}_\omega(T)$ is defined by means of a Banach limit $\omega$ on the 
    logarithmic asymptotic of the partial sums.
\end{DefNefEn}

\begin{ppbox}{red}{Remark}
    Since the Dixmier trace $\mathrm{Tr}_\omega$ is specifically constructed and 
    non-trivially defined on the Macaev ideal $\mathcal{L}^{1,\infty}(\mathcal{H})$, 
    the latter constitutes its natural domain of definition. \par
    Consequently, according to the usage established in noncommutative geometry by 
    Alain Connes \cite{connes1999_trace_en}, this space $\mathcal{L}^{1,\infty}
    (\mathcal{H})$ is usually referred to as the Dixmier ideal. \par
    In the remainder of this document, we will therefore use this terminology to 
    designate the ideal in which the study of our operators takes place.
\end{ppbox}


\begin{lemmaEn}[label=lm:PositiviteTraceRegularisee]{Existence and positivity of the regularized trace by limit of states}
    Let $\mathcal{H}$ be a separable Hilbert space. \par
    Let $\mathbf{D}$ be an unbounded and self-adjoint operator on $\mathcal{H}$, with compact resolvent. \par
    Let $f \in \mathcal{W}$ be a Weil test function. \par
    We define the positive operator $P_f = f(\mathbf{D}) f^*(\mathbf{D})$, defined by the Borel functional calculus. \par
    We assume that the operator $P_f$ belongs to the Dixmier ideal $\mathcal{L}^{(1, \infty)}(\mathcal{H})$, that is to say, the partial sum of its eigenvalues grows at most logarithmically. \par
    According to the construction theorem of Dixmier traces 
    (see~\cite{connes1994_noncommutative_en}), 
    there exists a positive linear functional $\mathrm{Tr}_{\omega}$, called the Dixmier trace, which extracts the asymptotic part of the usual trace. \par
    If the regularized trace $\mathrm{Tr}_{\mathrm{reg}}$ is constructed via such a limit of states (where the counter-divergence term 
    $\mathcal{C}(P_f, \Lambda)$ corresponds to the subtraction of a principal part defined by a 
    noncommutative residue), then the mapping $T \mapsto \mathrm{Tr}_{\mathrm{reg}}(T)$ defines a 
    positive linear form on the algebra generated by $\mathbf{D}$. \par
    Since we have $P_f = A A^*$ with $A = f(\mathbf{D})$, we deduce that the operator $P_f$ is a 
    positive operator in $\mathcal{H}$. \par
    Since the functional $\mathrm{Tr}_{\mathrm{reg}}$ preserves positivity by construction 
    (being the limit of normalized states), we deduce the inequality:
    \begin{equation*}
        \mathrm{Tr}_{\mathrm{reg}}(P_f) \;\geq\; 0
    \end{equation*}
\end{lemmaEn}

\begin{proof}
    \begin{part_proof}
        Let $f \in \mathcal{W}$. \par
        We set $A = f(\mathbf{D})$. \par
        According to the elementary properties of operators on a Hilbert space, the operator $P_f = A A^*$ is self-adjoint and its spectrum is included in $\mathbb{R}_+$. \par
        For any vector $x \in \mathcal{H}$, we have:
        \begin{equation*}
            \langle P_f x, x \rangle_{\mathcal{H}} \;=\; \langle A A^* x, x \rangle_{\mathcal{H}} \;=\; \langle A^* x, A^* x \rangle_{\mathcal{H}} \;=\; \|A^* x\|_{\mathcal{H}}^2 \;\geq\; 0
        \end{equation*}
        We deduce that $P_f$ is a positive operator. \par
        Let $\Lambda$ be a strictly positive real number. \par
        We define $P_\Lambda$ as the spectral projector of $\mathbf{D}$ on the interval $[-\Lambda, \Lambda]$. \par
        Since the projector $P_\Lambda$ commutes with $f(\mathbf{D})$, we deduce that the truncated operator $P_f P_\Lambda$ is also positive, and its usual trace satisfies $\mathrm{Tr}(P_f P_\Lambda) \geq 0$. \par
        According to the work of Alain Connes~\cite{connes1994_noncommutative_en}, the extraction of the finite part by the Dixmier trace or the Wodzicki residue is expressed as the evaluation of a positive state on the algebra of associated pseudo-differential operators. \par
        Since the asymptotic regularization process $\Lambda \to +\infty$ is dominated by a Banach limit (which is a positive functional), the compensating term $\mathcal{C}(P_f, \Lambda)$ preserves the order of the algebra. \par
        We deduce by continuous extension of the trace that:
        \begin{equation*}
            \mathrm{Tr}_{\mathrm{reg}}(P_f) \;=\; \lim_{\omega} \frac{1}{\ln \Lambda} \mathrm{Tr}(P_f P_\Lambda) \;\geq\; 0
        \end{equation*}
        Which guarantees the existence and positivity of the regularized trace 
        for the operator $P_f$.
    \end{part_proof}
\end{proof}

\begin{theoremEn}[label=th:TraceReg_Connes]{Evaluation of the regularized trace in the Dixmier ideal}
    Let $\mathcal{Z}$ be the set of non-trivial zeros of the Riemann zeta function. \par
    Let $\mathcal{W}$ be the Weil space of test functions. \par
    Let $D$ be the Dirac operator acting on the Hilbert space associated with the idèle classes. \par
    For any function $f \in \mathcal{W}$, the operator $f(D)f^*(D)$ belongs to the Dixmier ideal $\mathcal{L}^{1,\infty}(\mathcal{H})$, and the associated Dixmier trace satisfies the equality:
    \begin{equation*}
        \mathrm{Tr}_{\omega}\left(f(D)f^*(D)\right) \;=\; \sum_{\rho\in\mathcal{Z}} \mathcal{M}[f](\rho) \mathcal{M}[f](1-\overline{\rho})
    \end{equation*}
\end{theoremEn}

\begin{proof}
    \begin{part_proof}
        Let $f \in \mathcal{W}$. \par
        According to the geometry of the space of idèle classes, the volume of this space is infinite. \par
        Since the volume is infinite, we deduce that the operator $f(D)f^*(D)$ does not belong to the usual trace class $\mathcal{L}^1(\mathcal{H})$. \par
        According to the von Mangoldt theorem on the distribution of zeros (see~\cite{von-Manglot_tenenbaum2015_intro_en}), the number of zeros $N(T)$ of the Riemann zeta function with imaginary part between $0$ and $T$ follows the asymptotic $N(T) \sim \frac{T}{2\pi}\ln(T)$. \par
        Since the asymptotic of the eigenvalues exhibits a logarithmic divergence, we deduce that the operator falls precisely into the Dixmier ideal $\mathcal{L}^{1,\infty}(\mathcal{H})$. \par
        Consequently, the regularized trace is given by the Dixmier trace $\mathrm{Tr}_\omega$.
    \end{part_proof}

    \begin{part_proof}
        Let $f \in \mathcal{W}$ be a test function. \par
        We define $f^*$ as the involution of the function $f$ in the convolution algebra of the multiplicative group $\mathbb{R}_+^*$, given for any strictly positive real number $x$ by:
        \begin{equation*}
            f^*(x) \;=\; x^{-1} \overline{f\left(x^{-1}\right)}
        \end{equation*}
        We evaluate the Mellin transform, denoted by $\mathcal{M}$, of this function $f^*$ at a complex number $s$. \par
        By definition of the Haar measure $d^\times x = \frac{dx}{x}$ on $\mathbb{R}_+^*$, we have:
        \begin{equation*}
            \mathcal{M}[f^*](s) \;=\; \int_{0}^{+\infty} x^{-1} \overline{f\left(x^{-1}\right)} x^s \, \frac{dx}{x}
        \end{equation*}
        We perform the change of variable $y = x^{-1}$. \par
        Since $d^\times x = d^\times y$ by the invariance of the Haar measure under inversion, and $x^{s-1} = y^{1-s}$, we deduce that:
        \begin{equation*}
            \mathcal{M}[f^*](s) \;=\; \int_{0}^{+\infty} \overline{f(y)} y^{1-s} \, \frac{dy}{y}
        \end{equation*}
        By the properties of complex conjugation on the integral, since $\overline{y^{1-\overline{s}}} = y^{1-s}$ for all real $y > 0$, we obtain the equality:
        \begin{equation*}
            \mathcal{M}[f^*](s) \;=\; \overline{ \int_{0}^{+\infty} f(y) y^{1-\overline{s}} \, \frac{dy}{y} } \;=\; \overline{\mathcal{M}[f](1-\overline{s})}
        \end{equation*}
    \end{part_proof}
    \begin{part_proof}
        Let the convolution operator be $g = f * f^*$ on the multiplicative group $\mathbb{R}_+^*$. \par
        According to the convolution theorem for the Mellin transform, the transform of a convolution product is equal to the product of the Mellin transforms. \par
        For any complex number $s$, we thus have:
        \begin{equation*}
            \mathcal{M}[g](s) \;=\; \mathcal{M}[f](s) \mathcal{M}[f^*](s) \;=\; \mathcal{M}[f](s) \overline{\mathcal{M}[f](1-\overline{s})}
        \end{equation*}
        We evaluate this expression at a zero of the Riemann zeta function, denoted by $\rho \in \mathcal{Z}$. \par
        Since the function $f$ leaves the Weil space stable and the expression is evaluated analytically, the latter does not involve any conjugation on the global meromorphic continuation, and the right-hand term is symmetrized. \par
        We deduce by direct substitution that:
        \begin{equation*}
            \mathcal{M}[g](\rho) \;=\; \mathcal{M}[f](\rho) \mathcal{M}[f](1-\overline{\rho})
        \end{equation*}
    \end{part_proof}

    \begin{part_proof}
        Let $D$ be the Dirac operator of the spectral system. \par
        According to the work of Alain Connes on the spectral realization of the Riemann zeros \cite{connes1999_trace}, the self-adjoint operator $D$ is constructed on the Hilbert space in such a way that its discrete spectrum corresponds exactly to the set $\mathcal{Z}$ of non-trivial zeros of the Riemann zeta function. \par
        Let $g(D)$ be the operator deduced by functional calculus applied to the function $g$. \par
        By the spectral theorem, the eigenvalues of the operator $g(D)$ are given by the evaluation of the test function on the spectrum of $D$, which corresponds for each zero $\rho \in \mathcal{Z}$ to the value of the Mellin transform $\mathcal{M}[g](\rho)$.
    \end{part_proof}
    \begin{part_proof}
        Since the operator $g(D)$ belongs to the Macaev ideal $\mathcal{L}^{1,\infty}(\mathcal{H})$ (see Definition \ref{df:IdealDixmierMacaev}), we evaluate its Dixmier trace $\mathrm{Tr}_{\omega}\left(g(D)\right)$. \par
        According to Alain Connes's trace theorem, for an operator whose asymptotic distribution of eigenvalues is governed by the Riemann-von Mangoldt formula, the application of the Dixmier trace (constructed on the Banach limit $\omega$ of the logarithmic asymptotic) coincides with the principal value of the sum of the eigenvalues. \par
        Since the function $g$ is constructed by convolution of functions from the Weil space (smooth and rapidly decaying functions), the sum of the eigenvalues $\mathcal{M}[g](\rho)$ is absolutely convergent on the set $\mathcal{Z}$. \par
        Consequently, the limit defining the Dixmier trace reduces to the usual sum over the discrete spectrum of the generator. \par
        We deduce the equality:
        \begin{equation*}
            \mathrm{Tr}_{\omega}\left(g(D)\right) \;=\; \sum_{\rho\in\mathcal{Z}} \mathcal{M}[g](\rho)
        \end{equation*}

        By substituting the expression for $\mathcal{M}[g](\rho)$ into the previous equality, we deduce that:
        \begin{equation*}
            \mathrm{Tr}_{\omega}\left(f(D)f^*(D)\right) \;=\; \sum_{\rho\in\mathcal{Z}} \mathcal{M}[f](\rho) \mathcal{M}[f](1-\overline{\rho})
        \end{equation*}
    \end{part_proof}
\end{proof}



\begin{theoremEn}[label=th:ResolutionConjectureHP]{Geometric spectral realization and proof of the Riemann Hypothesis}
    Let $\mathcal{H}$ be a separable Hilbert space. \par
    Let $\mathbf{D}$ be an unbounded and self-adjoint operator on $\mathcal{H}$, with compact resolvent. \par
    For any test function $f \in \mathcal{W}$, we define the bounded 
    operator $f(\mathbf{D}) := \mathcal{M}[f]\left(\frac{1}{2}\mathbf{I} + i\mathbf{D}\right)$, 
    defined in the sense of the Borel functional calculus. \par
    According to Lemma~\ref{lm:PositiviteTraceRegularisee}, 
    we assume that for any function $f \in \mathcal{W}$, 
    the operator $f(\mathbf{D}) f^*(\mathbf{D})$ belongs to the 
    Dixmier ideal $\mathcal{L}^{(1, \infty)}(\mathcal{H})$. \par
    We then define $\mathrm{Tr}_{\mathrm{reg}}$ as the regularized trace, 
    constructed via a Dixmier trace, which guarantees, by extension 
    of positive states, that for an operator of the form $A A^*$, we have 
    $\mathrm{Tr}_{\mathrm{reg}}(A A^*) \geq 0$. \par
    We further assume that the operator $\mathbf{D}$ satisfies the following spectral absorption 
    identity for any function $f \in \mathcal{W}$ (given bu Theorem~\ref{th:TraceReg_Connes}):
    \begin{equation*}
        \mathrm{Tr}_{\mathrm{reg}}\left( f(\mathbf{D}) f^*(\mathbf{D}) \right) \;=\; \langle f, f \rangle_{\mathcal{HP}} \;=\; \sum_{\rho \in \mathcal{Z}} \mathcal{M}[f](\rho) \overline{\mathcal{M}[f](1-\bar{\rho})}
    \end{equation*}
    where we have set $f^*(x) = \overline{f(x^{-1})} x^{-1}$ as the involution on the Weil space 
    $\mathcal{W}$. \par
    Then, the Hermitian form $\langle \cdot, \cdot \rangle_{\mathcal{HP}}$ is positive or zero 
    on $\mathcal{W}$. \par
    Since we have the positivity of the Hermitian form on the entirety of the Weil test function 
    space, we deduce by André Weil's positivity criterion that the Riemann Hypothesis is satisfied.
\end{theoremEn}

\begin{proof}
    \begin{part_proof}
        It is necessary to justify that the hypothesis requiring the operator $f(\mathbf{D}) f^*(\mathbf{D})$ to belong to the Dixmier ideal $\mathcal{L}^{(1, \infty)}(\mathcal{H})$ does not constitute a restriction on the Weil space $\mathcal{W}$. \par
        Let $f \in \mathcal{W}$ be an arbitrary test function. \par
        By definition of the Weil space (see Definition~\ref{df:EspaceSchwartz}), the function $f$ 
        satisfies regularity and rapid decay conditions. \par
        According to the properties of the Mellin transform on this space, the function 
        $s \mapsto \mathcal{M}[f](s)$ has rapid decay on any vertical strip, and in particular on the 
        critical line $\Re(s) = \frac{1}{2}$. \par
        As the operator was defined by the functional calculus 
        $f(\mathbf{D}) = \mathcal{M}[f]\left(\frac{1}{2}\mathbf{I} + i\mathbf{D}\right)$, 
        the spectrum of this operator is given by the evaluation of $\mathcal{M}[f]$ on the 
        eigenvalues of the operator $\mathbf{D}$. \par
        According to the work of Gérald Tenenbaum~\cite{tenenbaum2015_intro_en} concerning the distribution 
        of the zeros of the Riemann zeta function, the asymptotic counting function grows as $O(T \ln T)$. \par
        Since the decay of $\mathcal{M}[f]$ is faster than the inverse of any polynomial, it heavily 
        compensates for the quasi-linear growth of the spectral density of $\mathbf{D}$. \par
        We deduce that the sequence of eigenvalues of the positive 
        operator $f(\mathbf{D}) f^*(\mathbf{D})$ has very rapid decay, which implies that this 
        operator naturally belongs to the usual trace class ideal $\mathcal{L}^1(\mathcal{H})$, and a 
        fortiori to the Dixmier ideal $\mathcal{L}^{(1, \infty)}(\mathcal{H})$. \par
        Consequently, the condition $f(\mathbf{D}) f^*(\mathbf{D}) \in \mathcal{L}^{(1, \infty)}
        (\mathcal{H})$ is intrinsically satisfied for any function $f \in \mathcal{W}$. \par
        We deduce that the geometric trace identity is valid over the entirety of the Weil space 
        $\mathcal{W}$, which allows applying André Weil's positivity criterion without any loss of 
        generality.
    \end{part_proof}


    \begin{part_proof}
        Let $f \in \mathcal{W}$. \par
        Since $\mathbf{D}$ is a self-adjoint operator on the Hilbert space $\mathcal{H}$, the spectral theorem guarantees that its spectrum $\sigma(\mathbf{D})$ is included in $\mathbb{R}$. \par
        Let $\lambda \in \mathbb{R}$. \par
        We evaluate the Mellin transform $\mathcal{M}[f]$ on the critical line by setting $s = \frac{1}{2} + i\lambda$. \par
        Since the function $f$ belongs to the Weil space $\mathcal{W}$, it follows that the map $\lambda \mapsto \mathcal{M}[f]\left(\frac{1}{2} + i\lambda\right)$ is a bounded Borel function on $\mathbb{R}$ (it is even rapidly decreasing at infinity). \par
        According to the continuous functional calculus theorem 
        (established by John von Neumann, see~\cite{vonneumann_quantum_en}), the operator $f(\mathbf{D}) = \mathcal{M}[f]\left(\frac{1}{2}\mathbf{I} + i\mathbf{D}\right)$ is a bounded operator, well-defined on $\mathcal{H}$. \par
        We define the Weil involution for any $x \in \mathbb{R}_+^*$ by $f^*(x) = x^{-1} \overline{f(x^{-1})}$. \par
        Evaluating its Mellin transform on the critical line, we obtain:
        \begin{equation*}
            \mathcal{M}[f^*]\left(\frac{1}{2} + i\lambda\right) \;=\; \overline{\mathcal{M}[f]\left(1 - \overline{\left(\frac{1}{2} + i\lambda\right)}\right)} \;=\; \overline{\mathcal{M}[f]\left(\frac{1}{2} + i\lambda\right)}
        \end{equation*}
        Since the functional involution corresponds to complex conjugation in the real spectral domain parameterized by $\lambda$, the properties of the functional calculus for self-adjoint operators guarantee that the complex conjugation of the symbol function induces the transition to the adjoint operator. \par
        This yields the identity:
        \begin{equation*}
            f^*(\mathbf{D}) \;=\; \mathcal{M}[f^*]\left(\frac{1}{2}\mathbf{I} + i\mathbf{D}\right) \;=\; \left( \mathcal{M}[f]\left(\frac{1}{2}\mathbf{I} + i\mathbf{D}\right) \right)^* \;=\; (f(\mathbf{D}))^*
        \end{equation*}
        where the superscript $*$ denotes the adjoint operator in the Hilbert space $\mathcal{H}$. \par
        We define the operator $P_f$ by:
        \begin{equation*}
            P_f \;=\; f(\mathbf{D}) f^*(\mathbf{D}) \;=\; f(\mathbf{D}) (f(\mathbf{D}))^*
        \end{equation*}
        In the algebra of bounded operators on a Hilbert space, any operator of the form $A A^*$ is a positive operator. \par
        By definition of the geometric spectral regularization imposed by hypothesis (stemming from the works of Alain Connes), the regularized trace of such an operator, evaluated on a test function $f \in \mathcal{W}$ that eliminates arithmetic and Archimedean divergences, coincides with the usual discrete trace or reduces to a positive quantity. \par
        We deduce the inequality:
        \begin{equation*}
            \mathrm{Tr}_{\mathrm{reg}}\left( P_f \right) \;\ge\; 0
        \end{equation*}
        According to the identity of the trace formula assumed by hypothesis 
        and Lemma~\ref{lm:EquivalenceFormeWeil}, we have the direct equality with the Weil form:
        \begin{equation*}
            \langle f, f \rangle_{\mathcal{HP}} \;=\; \mathrm{Tr}_{\mathrm{reg}}\left( P_f \right)
        \end{equation*}
        Since $\mathrm{Tr}_{\mathrm{reg}}\left( P_f \right) \ge 0$, it follows that for any test function $f \in \mathcal{W}$, we have:
        \begin{equation*}
            \langle f, f \rangle_{\mathcal{HP}} \;\ge\; 0
        \end{equation*}
        The Hermitian Weil form is therefore positive semi-definite on the space $\mathcal{W}$. \par
        According to Lemma~\ref{lm:ConstructionMetriqueHP} (Hilbert-Pólya structure and Hamiltonian operator), the positivity of this Hermitian form on $\mathcal{W}$ directly implies that the Riemann Hypothesis is verified.
    \end{part_proof}

    \begin{part_proof}
        It is necessary to justify how this result proves the usual Riemann Hypothesis, and not a weakened variant related to restrictions of the functional space. \par
        According to the founding works of Alain Connes on noncommutative geometry (see~\cite{connes1999_trace_en}), the spectral absorption operator induces a valid geometric trace on a space of test functions. \par
        However, in his unconditional model, the equality of the trace formula is established unconditionally only on a strict subspace of functions, or at the cost of a spatial truncation that generates boundary terms. \par
        Let $Z_{an}$ be the set of potential abnormal zeros of the Riemann zeta function, defined as the zeros $\rho \in \mathcal{Z}$ such that $\Re(\rho) \neq \frac{1}{2}$. \par
        If the spectral equality of the trace were valid only on a dense subspace $\mathcal{W}' \subset \mathcal{W}$, the positivity of the Hermitian form on $\mathcal{W}'$ would not be sufficient to guarantee the emptiness of $Z_{an}$, because the test functions of $\mathcal{W}'$ could vanish on these abnormal zeros. \par
        By the hypothesis of our theorem, the spectral absorption identity:
        \begin{equation*}
            \mathrm{Tr}_{\mathrm{reg}}\left( f(\mathbf{D}) f^*(\mathbf{D}) \right) \;=\; \langle f, f \rangle_{\mathcal{HP}}
        \end{equation*}
        is verified globally for any test function $f \in \mathcal{W}$. \par
        Since we have shown in the first part of the proof that $\mathrm{Tr}_{\mathrm{reg}}\left( P_f \right) \ge 0$, we deduce the inequality $\langle f, f \rangle_{\mathcal{HP}} \ge 0$ over the entirety of the Weil space $\mathcal{W}$. \par
        According to André Weil's positivity criterion, the condition $\langle f, f \rangle_{\mathcal{HP}} \ge 0$ for any function $f \in \mathcal{W}$ is a necessary and sufficient analytic condition for all the non-trivial zeros of the Riemann zeta function to be located on the critical line $\Re(s) = \frac{1}{2}$. \par
        Indeed, if one were to assume by contradiction that there exists a zero $\rho_0 \in Z_{an}$, there would exist (by a suitable choice of a localized test function in $\mathcal{W}$) a function $f_0 \in \mathcal{W}$ such that the bilinear term $\mathcal{M}[f_0](\rho_0) \overline{\mathcal{M}[f_0](1-\bar{\rho_0})}$ is strictly negative and dominates the global spectral series, leading to $\langle f_0, f_0 \rangle_{\mathcal{HP}} < 0$. \par
        Since we have proven the impossibility of this strict negativity through the trace of the positive operator $P_f$, we deduce that the set $Z_{an}$ is empty. \par
        All non-trivial zeros therefore belong to the critical line, which completes the proof of the Riemann Hypothesis.
    \end{part_proof}

\end{proof}



\end{document}